% Generated draft: requires mathematical and source review.
\documentclass[11pt]{article}
\usepackage[margin=25mm]{geometry}
\usepackage{fontspec}
\setmainfont{Latin Modern Roman}
\usepackage{amsmath,amssymb,mathtools,unicode-math}
\setmathfont{Latin Modern Math}
\usepackage{enumitem,needspace,xurl,hyperref}
\setlength{\parindent}{0pt}
\setlength{\parskip}{5pt}
\setlength{\emergencystretch}{3em}
\newcommand{\N}{\mathbb{N}}
\newcommand{\Z}{\mathbb{Z}}
\newcommand{\Q}{\mathbb{Q}}
\newcommand{\R}{\mathbb{R}}
\newcommand{\C}{\mathbb{C}}
\newcommand{\F}{\mathbb{F}}
\newcommand{\A}{\mathbb{A}}
\newcommand{\PP}{\mathbb{P}}
\newcommand{\E}{\mathbb{E}}
\newcommand{\eps}{\varepsilon}
\newcommand{\dd}{\,\mathrm{d}}
\begin{document}
\section*{UnsolvedMath ID: 30001348}
\textbf{Source problem code:} catalogue entry
\subsection{Definitions and mathematical statement}
Let \(A\) be a noetherian regular local \(k\)-algebra, with maximal ideal \(\mathfrak m\) and residue field \(\kappa=A/\mathfrak m\), and let \(\mathfrak a\subseteq A\) be an ideal. Regularity means that \(\dim_\kappa(\mathfrak m/\mathfrak m^2)=\dim A\). Retaining the source notation, put
\[
C_m(\mathfrak a):=\{f\in A:X_m(0)\subset H_m(0)\},\qquad
\mathfrak a^\ast:=\bigcap_{m\in\mathbb N}C_m(\mathfrak a).
\]
The supplied question is whether
\[
\forall A\ \forall\mathfrak a\subseteq A,\qquad \mathfrak a^\ast=\mathfrak a,
\]
where \(A\) ranges over the noetherian regular local \(k\)-algebras under consideration and \(\mathfrak a\) ranges over all their ideals.

\textbf{Unresolved formulation qualifications.} The supplied record does not define \(X\), \(H\), the constructions \(X_m(0)\) and \(H_m(0)\), or the dependence of \(H\) on \(f\). A possible interpretation is that \(X=\operatorname{Spec}(A/\mathfrak a)\), \(H=\operatorname{Spec}(A/(f))\), and the displayed objects are fibers of jet schemes over the closed point, but this interpretation is not specified by the supplied text and is not adopted as a definitive formulation here. In particular, containment as closed subschemes and containment of underlying sets give different conditions and must be distinguished. The role of \(0\), the convention for \(\mathbb N\), and any hypotheses on \(k\) also remain unspecified. Thus the displayed question preserves the supplied target but does not yet define a unique closure operation.
\subsection{Short English statement}
The question asks whether the stated jet-based closure fixes every ideal in every noetherian regular local \(k\)-algebra. The supplied record leaves the jet constructions and the meaning of containment undefined, so a unique precise formulation requires source clarification.
\subsection{Sources}
\begin{itemize}
\item Local UnsolvedMath registry, permanent ID \(30001348\), source problem code \texttt{catalogue entry}, title ``Jet-Closure Equality in Regular Local Algebras.'' The supplied record attributes the problem to the Oberwolfach Reports collection; no individual proposer is supplied.
\item Supplied original citation: \emph{Singularities (2010)}, Oberwolfach Reports, pp.~2405--2470, DOI \texttt{10.4171/owr/2009/43}. \url{https://doi.org/10.4171/owr/2009/43}. The registry separately supplies proposed year \(2009\).
\item Supplied literature reference, labeled ``jet closures and local isomorphism'': \url{https://doi.org/10.1016/j.jalgebra.2018.01.008}.
\item Supplied literature reference, labeled ``jet closures of singularities'': \url{https://arxiv.org/abs/2507.05796}.
\item The literature-sources list also includes ``original Oberwolfach report'': \url{https://doi.org/10.4171/owr/2009/43}.
\item Literature assessment dated exactly \texttt{2026-08-21}: the supplied curation labels the problem ``partially solved'' and states that jet and jet-support closures are related to integral closure and computed in important classes, while the universal equality asked in the report is not known. This is supplied dated curation derived from the OWR collection, not an independently verified literature assessment; the status and cited claims require independent verification.
\end{itemize}
\end{document}
